\honest{Making History Available}{Eliezer Yudkowsky}{2007}

I have a name for a habit of thought: the logical fallacy of generalization from fictional evidence. A journalist writing about AI brings up the \textsc{Terminator} films, not as prophecy, but as if they were a case that had happened somewhere.

The inverse error is to be moved too little by history. Why? ``In our ancestral environment, there were no movies; what you saw with your own eyes was true. Is it any wonder'' that films move us too much? \nb{The post gives no evidence for this account.} History reaches us as ink on paper, and we file it with the novels.

When I saw that my mysterious answer repeated the vitalists' mistake, I realized that vitalism had happened to real people. Had I lived through past revolutions in science, I say, ``I probably would not have made the same mistake again.'' \nb{Lord Kelvin, the vitalist Yudkowsky quotes elsewhere, lived through a revolution in his own field, the science of heat, and stayed a vitalist.}

So I resolved to think as if everything in the history books had happened to me. ``Is there so much difference'' between seeing an event and reading about it, I ask, since both reach you by causal chains? \nb{The post's footnote names one difference: history books over-represent rulers, so one should remember ``a thousand peasants for every ruler.'' The paragraph does not take it up.}

The Earth became older. I remembered Rome rising and falling, and the modern world became more fragile. So many mistakes are repeated because we do not remember making them. Don't you remember how many times your biases have killed you? ``I remember; I wasn't there.''

So, I end, the next time you doubt that the future will be strange, remember being born in a hunter-gatherer tribe, believing you could fly by eating the right mushrooms, and then flying. Remember thinking slavery was right, and changing your mind. Remember that you did not guess how the world would change. \nb{The post opened by saying that vividly imagined cases, such as films, come to be treated as evidence. Its remedy is to imagine history vividly in the first person, with particulars the reader must invent. It does not say how to keep those particulars from counting as evidence.}
